% !TEX program = xelatex
%=====================================================================
%  ECE / EECS One-Page Resume  (XeLaTeX)
%  Style reference: MIT EECS / CMU ECE / Stanford CS / Berkeley EECS
%  Compile with XeLaTeX (VSCode + LaTeX Workshop recipe: "xelatex")
%=====================================================================
\documentclass[10pt,a4paper]{article}

\usepackage[left=1.4cm,right=1.4cm,top=1.3cm,bottom=1.3cm]{geometry}
\usepackage{enumitem}
\usepackage{titlesec}
% \usepackage{titlerule}   % 若无此包，删除本行即可，下面的规则线仍可用
\usepackage{hyperref}
\usepackage{xcolor}
\usepackage[overload]{textcase} % LaTeX-aware \MakeUppercase，不破坏 \color 等命令
\usepackage{fontawesome5} % 图标：邮箱/GitHub 等（可选，缺失可注释）

% ---- 字体（XeLaTeX 专用，TeXLive 自带，无需下载）---------------------
\usepackage{fontspec}
\setmainfont{TeX Gyre Termes}      % 正文衬线体（类 Times，但更精致）
\newfontfamily\headingfont{TeX Gyre Heros} % 标题无衬线体（类 Helvetica）
% 中文字体支持
\usepackage{xeCJK}
\setCJKmainfont{SimSun}[AutoFakeBold=2.5, AutoFakeSlant=0.2]   % 中文正文（宋体），自动伪粗体/伪斜体
\setCJKsansfont{Microsoft YaHei}[AutoFakeBold=2.5]              % 中文无衬线（微软雅黑），自动伪粗体
% 若编译报字体错误，注释掉上面几行 fontspec/xeCJK 相关设置即可。
% ---------------------------------------------------------------------

\pagestyle{empty}

\hypersetup{
  colorlinks=true,
  urlcolor=black,
  linkcolor=black
}

% ---- 段落列表样式 ----------------------------------------------------
\setlist[itemize]{
  leftmargin=1.3em,
  itemsep=1.5pt,
  topsep=2pt,
  parsep=0pt,
  label=\textbullet
}

% ---- 主色调 ---------------------------------------------------------
\definecolor{mainblue}{HTML}{00509E}
\colorlet{MAINBLUE}{mainblue}

% ---- Section 标题：全大写 + 下方整条横线 ------------------------------
\titleformat{\section}
  {\headingfont\large\bfseries\color{mainblue}\MakeUppercase}
  {}{0em}{}[\vspace{-6pt}{\color{mainblue}\rule{\textwidth}{1.2pt}}]
\titlespacing*{\section}{0pt}{10pt}{6pt}

% ---- 便捷宏 ----------------------------------------------------------
% 项目条目：左标题（加粗），右侧可放技术栈/时间
\newcommand{\project}[2]{%
  \vspace{3pt}%
  \noindent\textbf{#1}\hfill{\small\textit{#2}}\par
  \vspace{1pt}%
}

% 学历/经历条目：左机构，右时间
\newcommand{\entry}[2]{%
  \noindent\textbf{#1}\hfill{#2}\par
}

\begin{document}

%=====================================================================
%  HEADER  —— 姓名与联系方式（请填写）
%=====================================================================
\begin{center}
  {\headingfont\Huge\bfseries your name}\\[4pt]
  {\normalsize ECE Undergraduate | Hardware + Software Builder}\\[6pt]
  {\small
    \faEnvelope\ \href{mailto:your mail address}{your mail address}
    \quad|\quad
    \faPhone\ +86 your phone number
    \quad|\quad
    \faGithub\ \href{https://github.com/yourname}{github.com/yourname}
  }
\end{center}

%=====================================================================
%  EDUCATION
%=====================================================================
\section{Education}

\entry{Zhejiang University}{Expected Graduation: 20xx}
Bachelor of Engineering in Electronic and Computer Engineering \hfill GPA: \textbf{4.00} / 4.00\par %这里请改成你的专业和成绩
\vspace{2pt}
% \textbf{Relevant Coursework:} Signals and Systems, Digital Circuits, Analog Circuits,
% Electronic circuit design experiment, Artificial Intelligence, Electromagnetic Field and Wave...

%=====================================================================
%  TECHNICAL SKILLS
%=====================================================================
\section{Technical Skills}

%这里提到的软件和技术只是举个例子，具体请根据你自己的情况修改

\begin{tabular}{@{}p{3.2cm}@{}p{14cm}@{}}
  \textbf{Programming}      & Python (Advanced), C/C++, MATLAB \\[3pt]
  % \textbf{AI / Vision}      & OpenCV, NumPy, PyTorch, scikit-learn \\[3pt]
  \textbf{Embedded}         & STM32, UART, Xilinx, Raspberry Pi \\[3pt]
  \textbf{Hardware}         & Altium Designer, Pspice, Oscilloscope, Signal Generator\\[3pt]
  \textbf{Tools}            & Git, VS Code, Keil, STM32CubeMX \\
\end{tabular}

%=====================================================================
%  SELECTED PROJECTS
%=====================================================================
\section{Builds \& Projects}

% ---- 1. sth.1 ----
\project{your project name}{subject / tech stack}
\begin{itemize}
  \item achieved something
  \item implemented something
\end{itemize}

% ---- 2. sth.2 ----
\project{your project name}{subject / tech stack}
\begin{itemize}
  \item achieved something
  \item implemented something
\end{itemize}

% ---- 3. sth.3 ----
\project{your project name}{subject / tech stack}
\begin{itemize}
  \item achieved something
  \item implemented something
\end{itemize}

% ---- 4. sth.4 ----
\project{your project name}{subject / tech stack}
\begin{itemize}
  \item achieved something
  \item implemented something
\end{itemize}


\section{Research Experience}

\entry{your institution}{time period}
\textit{Advisor: Professor xx}\
\begin{itemize}
  \item what you did in the research
  \item what you learned.
\end{itemize}

\entry{your institution}{time period}
\textit{Advisor: Professor xx}\
\begin{itemize}
  \item what you did in the research
  \item what you learned.
\end{itemize}
%=====================================================================
%  HONORS & AWARDS  (可选，暂留空)
%=====================================================================
\section{Honors \& Awards}
\begin{itemize}
  \item sth. you won
  \item sth. you won
\end{itemize}
%=====================================================================
%  LANGUAGES
%=====================================================================
\section{Languages}
Chinese (Native) \quad|\quad English (\textbf{CET-4}: 665 / \textbf{CET-6}: 621)

%=====================================================================
%  PAGE 2: CHINESE VERSION
%=====================================================================
\newpage

% 中文版重新定义 section 标题：取消 MakeUppercase，保留相同横线样式
\titleformat{\section}
  {\headingfont\large\bfseries\color{mainblue}}
  {}{0em}{}[\vspace{-6pt}{\color{mainblue}\rule{\textwidth}{1.2pt}}]

\begin{center}
  {\headingfont\Huge\bfseries 你的名字}\\[4pt]
  {\normalsize 某专业本科在读 | 软硬件协同开发}\\[6pt]
  {\small
    \faEnvelope\ \href{mailto:邮箱}{邮箱}
    \quad|\quad
    \faPhone\ +86 电话
    \quad|\quad
    \faGithub\ \href{https://github.com/}{github}
  }
\end{center}

\section{教育背景}

\entry{浙江大学}{预计毕业：20xx}
某专业工学学士 \hfill 平均绩点：\textbf{4.00} / 4.00\par
\textbf{核心课程}： 。。。等

\section{技术能力}

\begin{tabular}{@{}p{3.2cm}@{}p{14cm}@{}}
  \textbf{编程语言}   & Python、C/C++、MATLAB \\[3pt]
  \textbf{嵌入式开发}     & STM32、UART、Xilinx、Raspberry Pi \\[3pt]
  \textbf{硬件开发}       & Altium Designer、Pspice等 \\[3pt]
  \textbf{工具使用}       & Git、VS Code、Keil、STM32CubeMX \\
\end{tabular}

\section{项目经历}

\project{\textbf{。。。}}{。。。}
\begin{itemize}
  \item 开发了。。。
  \item 实现了。。。
\end{itemize}

\project{\textbf{。。。}}{。。。}
\begin{itemize}
  \item 构建了。。。
  \item 完成了。。。
  \item 。。。
\end{itemize}

\project{\textbf{。。。}}{。。。}
\begin{itemize}
  \item 构建了。。。
  \item 完成了。。。
  \item 。。。
\end{itemize}

\project{\textbf{。。。}}{。。。}
\begin{itemize}
  \item 构建了。。。
  \item 完成了。。。
  \item 。。。
\end{itemize}

\section{科研经历}

\entry{\textbf{浙江大学。。。学院 }}{时间}
\textit{导师：xx 教授}\
\begin{itemize}
  \item 参与。。。
  \item 。。。
\end{itemize}

\entry{\textbf{浙江大学。。。学院 }}{时间}
\textit{导师：xx 教授}\
\begin{itemize}
  \item 参与。。。
  \item 。。。
\end{itemize}

\section{荣誉与奖项}
\begin{itemize}
  \item 。。
  \item 。。
  \item 。。
\end{itemize}

\section{语言能力}
中文（母语）\quad|\quad 英语（\textbf{英语四级}：665 / \textbf{英语六级}：621）

\end{document}